%%%%%%%%%%%%%%%%
%%% exod 15 %%%
%%%%%%%%%%%%%%%%

\Chapter{15}

\Verse{1}
{
	\hJ{
	אָז יָשִׁיר מֹשֶׁה וּבְנֵי יִשְׂרָאֵל אֶת הַשִּׁירָה הַזֹּאת לַיהֹוָה וַיֹּאמְרוּ לֵאמֹר		
	}
	\hProto{
אָשִׁירָה לַיהֹוָה כִּי גָאֹה גָּאָה סוּס וְרֹכְבוֹ רָמָה בַיָּם׃			
	}
}
{
	\eJ{
		Then Moses and the children of Israel sang this song to
		YHWH. And they said, saying:
	}
	\eProto{
		Let me sing to YHWH, for He triumphed! Horse and its rider He cast in the sea.
	}
}
{
%	\footnote{
%		rider. This refers to those who are riding in the chariots. Cavalry was
%		not used in Egypt in this period.
%	}
}

\Verse{2}
{
	\hProto{עזִּי וְזִמְרָת יָהּ וַיְהִי לִי לִישׁוּעָה זֶה אֵלִי וְאַנְוֵהוּ אֱלֹהֵי אָבִי וַאֲרֹמְמֶנְהוּ׃}
}
{
	\eProto{
		My strength and song are Yah, and He became a salvation
		for me. This is my {\God}, and I’ll praise Him, my father’s
		{\God}, and I’ll hail Him.
	}
}
{
}

\Verse{3}
{
	\hProto{יְהֹוָה אִישׁ מִלְחָמָה יְהֹוָה שְׁמוֹ׃}
}
{
	\eProto{YHWH is a warrior. YHWH is His name.}
}
{
}

\Verse{4}
{
	\hProto{מַרְכְּבֹת פַּרְעֹה וְחֵילוֹ יָרָה בַיָּם וּמִבְחַר שָׁלִשָׁיו טֻבְּעוּ בְיַם סוּף׃}
}
{
	\eProto{
		Pharaoh’s chariots and his army He plunged in the sea and
		the choice of his troops drowned in the Red Sea.
	}
}
{
}

\Verse{5}
{
	\hProto{תְּהֹמֹת יְכַסְיֻמוּ יָרְדוּ בִמְצוֹלֹת כְּמוֹ אָבֶן׃}
}
{
	\eProto{
		The deeps covered them. They sank in the depths like a
		stone.
	}
}
{
}

\Verse{6}
{
	\hProto{יְמִינְךָ יְהֹוָה נֶאְדָּרִי בַּכֹּחַ יְמִינְךָ יְהֹוָה תִּרְעַץ אוֹיֵב׃}
}
{
	\eProto{
		Your right hand, YHWH, awesome in power, your right hand,
		YHWH, crushed the foe.
	}
}
{
}

\Verse{7}
{
	\hProto{וּבְרֹב גְּאוֹנְךָ תַּהֲרֹס קָמֶיךָ תְּשַׁלַּח חֲרֹנְךָ יֹאכְלֵמוֹ כַּקַּשׁ׃}
}
{
	\eProto{
		And in your triumph’s greatness you threw down your
		adversaries. You let go your fury: it consumed them like
		straw.
	}
}
{
}

\Verse{8}
{
	\hProto{וּבְרוּחַ אַפֶּיךָ נֶעֶרְמוּ מַיִם נִצְּבוּ כְמוֹ נֵד נֹזְלִים קָפְאוּ תְהֹמֹת בְּלֶב יָם׃}
}
{
	\eProto{
		And by wind from your nostrils water was massed, surf
		piled up like a heap, the deeps congealed in the heart of
		the sea.
	}
}
{
}

\Verse{9}
{
	\hProto{אָמַר אוֹיֵב אֶרְדֹּף אַשִּׂיג אֲחַלֵּק שָׁלָל תִּמְלָאֵמוֹ נַפְשִׁי אָרִיק חַרְבִּי תּוֹרִישֵׁמוֹ יָדִי׃}
}
{
	\eProto{
		The enemy said, “I’ll pursue! I’ll catch up! I’ll divide
		spoil! My soul will be sated! I’ll unsheathe my sword! My
		hand will deprive them!”
	}
}
{
}

\Verse{10}
{
	\hProto{נָשַׁפְתָּ בְרוּחֲךָ כִּסָּמוֹ יָם צָלְלוּ כַּעוֹפֶרֶת בְּמַיִם אַדִּירִים׃}
}
{
	\eProto{
		You blew with your wind. Sea covered them. They sank like
		lead in the awesome water.
	}
}
{
}

\Verse{11}
{
	\hProto{מִי כָמֹכָה בָּאֵלִם יְהֹוָה מִי כָּמֹכָה נֶאְדָּר בַּקֹּדֶשׁ נוֹרָא תְהִלֹּת עֹשֵׂה פֶלֶא׃}
}
{
	\eProto{
		Who is like you among the gods, YHWH! Who is like you:
		awesome in holiness! fearsome with splendors! making
		miracles!
	}
}
{
}

\Verse{12}
{
	\hProto{נָטִיתָ יְמִינְךָ תִּבְלָעֵמוֹ אָרֶץ׃}
}
{
	\eProto{You reached your right hand: earth swallowed them.}
}
{
}

\Verse{13}
{
	\hProto{נָחִיתָ בְחַסְדְּךָ עַם זוּ גָּאָלְתָּ נֵהַלְתָּ בְעזְּךָ אֶל נְוֵה קדְשֶׁךָ׃}
}
{
	\eProto{
		You led, in your kindness, the people you redeemed; you
		ushered, in your strength, to your holy abode.
	}
}
{
}

\Verse{14}
{
	\hProto{שָׁמְעוּ עַמִּים יִרְגָּזוּן חִיל אָחַז יֹשְׁבֵי פְּלָשֶׁת׃}
}
{
	\eProto{
		Peoples heard—they shuddered. Shaking seized Philistia’s
		residents.
	}
}
{
%	\footnote{
%		Peoples heard. This is not just about Egypt. As in the matter of the
%		plagues, the object is not only to defeat or punish the Egyptians. It is
%		to make YHWH known in the world. So the Song of the Sea declares that
%		others hear of this event: Philistia, Edom, Moab, all of Canaan’s
%		residents.
%	}
}

\Verse{15}
{
	\hProto{אָז נִבְהֲלוּ אַלּוּפֵי אֱדוֹם אֵילֵי מוֹאָב יֹאחֲזֵמוֹ רָעַד נָמֹגוּ כֹּל יֹשְׁבֵי כְנָעַן׃}
}
{
	\eProto{
		Then Edom’s chiefs were terrified. Moab’s chieftains:
		trembling seized them. All Canaan’s residents melted.
	}
}
{
}

\Verse{16}
{
	\hProto{תִּפֹּל עֲלֵיהֶם אֵימָתָה וָפַחַד בִּגְדֹל זְרוֹעֲךָ יִדְּמוּ כָּאָבֶן עַד יַעֲבֹר עַמְּךָ יְהֹוָה עַד יַעֲבֹר עַם זוּ קָנִיתָ׃}
}
{
	\eProto{
		Terror and fear came over them. At the power of your arm
		they’re silent like stone. ’Til your people passed, YHWH,
		’til the people you created passed.
	}
}
{
}

\Verse{17}
{
	\hProto{תְּבִאֵמוֹ וְתִטָּעֵמוֹ בְּהַר נַחֲלָתְךָ מָכוֹן לְשִׁבְתְּךָ פָּעַלְתָּ יְהֹוָה מִקְּדָשׁ אֲדֹנָי כּוֹנְנוּ יָדֶיךָ׃}
}
{
	\eProto{
		You’ll bring them and plant them in your legacy’s
		mountain, your throne’s platform, that you made, YHWH; a
		holy place, Lord, that your hands reared.
	}
}
{
}

\Verse{18}
{
	\hProto{יְהֹוָה יִמְלֹךְ לְעֹלָם וָעֶד׃}
}
{
	\eProto{YHWH will reign forever and ever!}
}
{
}

\Verse{19}
{
	\hR{כִּי בָא סוּס פַּרְעֹה בְּרִכְבּוֹ וּבְפָרָשָׁיו בַּיָּם וַיָּשֶׁב יְהֹוָה עֲלֵהֶם אֶת מֵי הַיָּם וּבְנֵי יִשְׂרָאֵל הָלְכוּ בַיַּבָּשָׁה בְּתוֹךְ הַיָּם׃}
}
{
	\eR{
		Because Pharaoh’s horses with his chariots and with his
		horsemen came in the sea, and YHWH brought back the water
		of the sea over them, and the children of Israel had gone
		on the dry ground through the sea.
	}
}
{
}

\Verse{20}
{
	\hProto{וַתִּקַּח מִרְיָם הַנְּבִיאָה אֲחוֹת אַהֲרֹן אֶת הַתֹּף בְּיָדָהּ וַתֵּצֶאןָ כל הַנָּשִׁים אַחֲרֶיהָ בְּתֻפִּים וּבִמְחֹלֹת׃}
}
{
	\eE{
		And Miriam, the prophetess, Aaron’s sister, took a drum in
		her hand, and all the women went out behind her with drums
		and with dances.
	}
}
{
%	\footnote{
%		Miriam, the prophetess. Here we are informed that she is a prophetess.
%		This will be crucial later, when she and Aaron will speak against Moses,
%		saying, “Has YHWH only just spoken through Moses? Hasn’t He also spoken
%		through us?” (Num 12:2).
%	}
}

\Verse{21}
{
	\hProto{וַתַּעַן לָהֶם מִרְיָם שִׁירוּ לַיהֹוָה כִּי גָאֹה גָּאָה סוּס וְרֹכְבוֹ רָמָה בַיָּם׃}
}
{
	\eE{
		And Miriam sang to them: Sing to YHWH for He triumphed!
		Horse and its rider He cast in the sea.
	}
}
{
%	\footnote{
%		Miriam sang. This song is therefore known both as the Song of the Sea
%		and as the Song of Miriam. My teacher Frank Moore Cross and my senior
%		colleague David Noel Freedman demonstrated together that this is one of
%		the oldest—if not the oldest—passages in the Bible. The Red Sea episode
%		was thus an early and decisive element of Israel’s story.
%	}
}

\Verse{22}
{
	\hJ{וַיַּסַּע מֹשֶׁה אֶת יִשְׂרָאֵל מִיַּם סוּף וַיֵּצְאוּ אֶל מִדְבַּר שׁוּר וַיֵּלְכוּ שְׁלֹשֶׁת יָמִים בַּמִּדְבָּר וְלֹא מָצְאוּ מָיִם׃}
}
{
	\eJ{
		And Moses had Israel travel from the Red Sea, and they
		went out to the wilderness of Shur. And they went three
		days in the wilderness. And they did not find water.
	}
}
{
%	\footnote{
%		they went three days in the wilderness. The original request they made
%		of Pharaoh—a three-day trip in the wilderness—is now fulfilled. And they
%		are not returning!
%	}
}

\Verse{23}
{
	\hJ{וַיָּבֹאוּ מָרָתָה וְלֹא יָכְלוּ לִשְׁתֹּת מַיִם מִמָּרָה כִּי מָרִים הֵם עַל כֵּן קָרָא שְׁמָהּ מָרָה׃}
}
{
	\eJ{
		And they came to Marah, and they were not able to drink
		water from Marah because it was bitter. On account of this
		he called its name Marah.
	}
}
{
}

\Verse{24}
{
	\hJ{וַיִּלֹּנוּ הָעָם עַל מֹשֶׁה לֵּאמֹר מַה נִּשְׁתֶּה׃}
}
{
	\eJ{
		And the people complained at Moses, saying, “What shall we
		drink?”
	}
}
{
%	\footnote{
%		the people complained at Moses. The narrative has proceeded from the
%		confrontation between Moses and the deity to that between Moses and the
%		Pharaoh, and now to the ongoing confrontation between Moses and the
%		people. And, as I indicated earlier, each new stage of confrontation
%		carries with it the background of what has come before, so that Moses’
%		experiences with the people—involving sometimes wisdom, sometimes rage,
%		and sometimes affection—must be read with some appreciation of the
%		growth that has taken place in his personality through his earlier
%		encounters with divinity and royalty. Having defeated the leader of
%		Egypt, Moses now finds himself to be the leader of a people, a people
%		who never chose him to lead them. Through a series of crises, the people
%		doubt his competence and rebel against his authority.
%	}
}

\Verse{25}
{
	\hJ{וַיִּצְעַק אֶל יְהֹוָה וַיּוֹרֵהוּ יְהֹוָה עֵץ וַיַּשְׁלֵךְ אֶל הַמַּיִם וַיִּמְתְּקוּ הַמָּיִם שָׁם שָׂם לוֹ חֹק וּמִשְׁפָּט וְשָׁם נִסָּהוּ׃}
}
{
	\eJ{
		And he cried to YHWH, and YHWH showed him a tree, and he
		threw it into the water, and the water sweetened. He set
		law and judgment for them there, and He tested them there.
	}
}
{
}

\Verse{26}
{
	\hOther{וַיֹּאמֶר אִם שָׁמוֹעַ תִּשְׁמַע לְקוֹל יְהֹוָה אֱלֹהֶיךָ וְהַיָּשָׁר בְּעֵינָיו תַּעֲשֶׂה וְהַאֲזַנְתָּ לְמִצְוֺתָיו וְשָׁמַרְתָּ כּל חֻקָּיו כּל הַמַּחֲלָה אֲשֶׁר שַׂמְתִּי בְמִצְרַיִם לֹא אָשִׂים עָלֶיךָ כִּי אֲנִי יְהֹוָה רֹפְאֶךָ׃}
}
{
	\eOther{
		And He said, “If you’ll listen to the voice of YHWH, your
		{\God}, and you’ll do what is right in His eyes and turn your
		ear to His commandments and observe all His laws, I won’t
		set on you any of the sickness that I set on Egypt,
		because I, YHWH, am your healer.”
	}
}
{
}

\Verse{27}
{
	\hP{וַיָּבֹאוּ אֵילִמָה וְשָׁם שְׁתֵּים עֶשְׂרֵה עֵינֹת מַיִם וְשִׁבְעִים תְּמָרִים וַיַּחֲנוּ שָׁם עַל הַמָּיִם׃}
}
{
	\eP{
		And they came to Elim. And twelve springs of water and
		seventy palm trees were there. And they camped there by
		the water.
	}
}
{
}
